% ---- score2 standard preamble (from: python3 score2_tools.py preamble) ----
\documentclass[10pt]{extarticle}
\usepackage[a4paper,margin=1.5cm]{geometry}
\usepackage[T1]{fontenc}
\usepackage{amsmath,amssymb}
\usepackage{graphicx}
\usepackage[export]{adjustbox}
\usepackage{multicol}
\usepackage{enumitem}
\usepackage{booktabs,array,tabularx}
\setlength{\columnsep}{0.7cm}
\setlength{\parindent}{0pt}
\setlength{\parskip}{3pt}
\pagestyle{empty}
\makeatletter
% figure registry (figdefs.tex is generated by the extract step)
\newcommand{\qfigdef}[4]{\@namedef{qf@file@#1}{#2}\@namedef{qf@w@#1}{#3}\@namedef{qf@wide@#1}{#4}}
\newcommand{\qf@check}[1]{\@ifundefined{qf@file@#1}{\PackageError{score2}{Unknown figure id #1}{}}{}}
\newcommand{\qf@img}[4][]{\includegraphics[width=#2\dimexpr\@nameuse{qf@w@#4}pt\relax,max width=#3,max height=0.45\textheight,#1]{\@nameuse{qf@file@#4}}}
% Full-width questions: finalize wraps every question that contains a wide figure in \qfullbegin...\qfullend,
% so the WHOLE question (stem, figures, options) is set across the page in reading order, then the two columns
% resume. The question's head - its number, stem and first figure (\qfullsplit marks the end of the head) -
% is kept on one page, so a number is never left at a page foot with its figure overleaf; the rest may break
% across pages normally, so no big gaps are left. Without a \qfullsplit, the whole question is kept together
% if it is under 45% of a page. Never write these yourself.
\newif\ifqf@full
\newif\ifqf@boxed
\newsavebox{\qf@box}
\newcommand{\qf@open}{\global\qf@fulltrue\par\global\qf@boxedtrue
  \setbox\qf@box\vbox\bgroup\hsize\textwidth\linewidth\textwidth}
\newcommand{\qf@emit}[1]{\par\egroup\global\qf@boxedfalse
  \ifdim\dimexpr\ht\qf@box+\dp\qf@box\relax<#1\textheight\noindent\usebox{\qf@box}\par\else\unvbox\qf@box\fi}
\newcommand{\qfullsplit}{\ifqf@boxed\qf@emit{0.85}\fi}
\newcommand{\qf@close}{\ifqf@boxed\qf@emit{0.45}\fi\global\qf@fullfalse}
\newcommand{\qfullbegin}{\end{multicols}\qf@open}
\newcommand{\qfullend}{\qf@close\begin{multicols}{2}}
\newcommand{\qfullstart}{\qf@open}
\newcommand{\qfulldone}{\qf@close}
% \qfig[scale]{id}: block figure, centred, directly where it is written (size set by the tool).
% Fallback if a wide figure is not inside a full-width question: break out of the columns in place - no floats.
\newcommand{\qfig}[2][1]{\qf@check{#2}%
  \ifqf@full\qf@blk{#1}{#2}\else
  \ifnum\@nameuse{qf@wide@#2}=1 \expandafter\@firstoftwo\else\expandafter\@secondoftwo\fi
  {\end{multicols}\begin{center}\qf@img{#1}{\textwidth}{#2}\end{center}\begin{multicols}{2}}%
  {\qf@blk{#1}{#2}}\fi}
\newcommand{\qf@blk}[2]{\par\smallskip{\centering\qf@img{#1}{\linewidth}{#2}\par}\smallskip}
% \qfigi[scale]{id}: inline figure (inside option lists / tables / minipages); never breaks columns.
\newcommand{\qfigi}[2][1]{\qf@check{#2}\qf@img[valign=c]{#1}{\linewidth}{#2}}
\makeatother
% structure macros - lint reads these to check completeness and figure placement
\newcommand{\papersec}[1]{\par\bigskip{\centering\Large\bfseries PAPER-#1\par}\medskip}
\newcommand{\partsec}[2]{\par\medskip\noindent{\large\bfseries PART-#1: #2}\par\smallskip}
\newcommand{\sect}[1]{\par\smallskip\noindent\textbf{#1}\par}
\newcommand{\qn}[1]{\par\smallskip\noindent\textbf{#1.}~}
\newcommand{\paraq}[2]{\par\smallskip\noindent\textbf{Paragraph for Questions #1 and #2}\par}
\newenvironment{qoptions}{\begin{itemize}[leftmargin=1.8em,itemsep=1pt,topsep=2pt,label={}]}{\end{itemize}}
\newcommand{\opt}[2]{\item[(#1)] #2}
% figure registry (generated by score2_tools.py finalize)
\qfigdef{p1_phy_s1_q09_1}{Score-II - Test 25 (20 Mar 2024)-img/p1_phy_s1_q09_1.png}{92.3}{0}
\qfigdef{p1_phy_s1_q11_5}{Score-II - Test 25 (20 Mar 2024)-img/p1_phy_s1_q11_5.png}{112.5}{0}
\qfigdef{p1_phy_s1_q12_1}{Score-II - Test 25 (20 Mar 2024)-img/p1_phy_s1_q12_1.png}{59.6}{0}
\qfigdef{p1_phy_s1_q15_1}{Score-II - Test 25 (20 Mar 2024)-img/p1_phy_s1_q15_1.png}{94.3}{0}
\qfigdef{p1_phy_s2_q01_1}{Score-II - Test 25 (20 Mar 2024)-img/p1_phy_s2_q01_1.png}{55.4}{0}
\qfigdef{p1_phy_s2_q03_1}{Score-II - Test 25 (20 Mar 2024)-img/p1_phy_s2_q03_1.png}{81.6}{0}
\qfigdef{p1_phy_s2_q09_1}{Score-II - Test 25 (20 Mar 2024)-img/p1_phy_s2_q09_1.png}{158.4}{0}
\qfigdef{p1_phy_s2_q10_1}{Score-II - Test 25 (20 Mar 2024)-img/p1_phy_s2_q10_1.png}{90.8}{0}
\qfigdef{p1_chem_s1_q01_1}{Score-II - Test 25 (20 Mar 2024)-img/p1_chem_s1_q01_1.png}{120.1}{0}
\qfigdef{p1_chem_s1_q02_1}{Score-II - Test 25 (20 Mar 2024)-img/p1_chem_s1_q02_1.png}{157.3}{0}
\qfigdef{p1_chem_s1_q02_2}{Score-II - Test 25 (20 Mar 2024)-img/p1_chem_s1_q02_2.png}{23.7}{0}
\qfigdef{p1_chem_s1_q02_3}{Score-II - Test 25 (20 Mar 2024)-img/p1_chem_s1_q02_3.png}{32.8}{0}
\qfigdef{p1_chem_s1_q02_4}{Score-II - Test 25 (20 Mar 2024)-img/p1_chem_s1_q02_4.png}{23.7}{0}
\qfigdef{p1_chem_s1_q02_5}{Score-II - Test 25 (20 Mar 2024)-img/p1_chem_s1_q02_5.png}{27.4}{0}
\qfigdef{p1_chem_s1_q03_1}{Score-II - Test 25 (20 Mar 2024)-img/p1_chem_s1_q03_1.png}{148.4}{0}
\qfigdef{p1_chem_s1_q04_1}{Score-II - Test 25 (20 Mar 2024)-img/p1_chem_s1_q04_1.png}{49.2}{0}
\qfigdef{p1_chem_s1_q04_2}{Score-II - Test 25 (20 Mar 2024)-img/p1_chem_s1_q04_2.png}{53.9}{0}
\qfigdef{p1_chem_s1_q05_1}{Score-II - Test 25 (20 Mar 2024)-img/p1_chem_s1_q05_1.png}{149.9}{0}
\qfigdef{p1_chem_s1_q05_3}{Score-II - Test 25 (20 Mar 2024)-img/p1_chem_s1_q05_3.png}{77.9}{0}
\qfigdef{p1_chem_s1_q05_2}{Score-II - Test 25 (20 Mar 2024)-img/p1_chem_s1_q05_2.png}{50.7}{0}
\qfigdef{p1_chem_s1_q06_1}{Score-II - Test 25 (20 Mar 2024)-img/p1_chem_s1_q06_1.png}{166.3}{0}
\qfigdef{p1_chem_s1_q07_1}{Score-II - Test 25 (20 Mar 2024)-img/p1_chem_s1_q07_1.png}{45.8}{0}
\qfigdef{p1_chem_s1_q07_2}{Score-II - Test 25 (20 Mar 2024)-img/p1_chem_s1_q07_2.png}{25.8}{0}
\qfigdef{p1_chem_s1_q08_1}{Score-II - Test 25 (20 Mar 2024)-img/p1_chem_s1_q08_1.png}{55.3}{0}
\qfigdef{p1_chem_s1_q08_2}{Score-II - Test 25 (20 Mar 2024)-img/p1_chem_s1_q08_2.png}{58.7}{0}
\qfigdef{p1_chem_s1_q08_3}{Score-II - Test 25 (20 Mar 2024)-img/p1_chem_s1_q08_3.png}{65.4}{0}
\qfigdef{p1_chem_s2_q01_1}{Score-II - Test 25 (20 Mar 2024)-img/p1_chem_s2_q01_1.png}{78.1}{0}
\qfigdef{p1_chem_s2_q02_1}{Score-II - Test 25 (20 Mar 2024)-img/p1_chem_s2_q02_1.png}{205.8}{0}
\qfigdef{p1_chem_s2_q03_1}{Score-II - Test 25 (20 Mar 2024)-img/p1_chem_s2_q03_1.png}{159.2}{0}
\qfigdef{p2_phy_s1_q01_1}{Score-II - Test 25 (20 Mar 2024)-img/p2_phy_s1_q01_1.png}{95.0}{0}
\qfigdef{p2_phy_s1_q02_1}{Score-II - Test 25 (20 Mar 2024)-img/p2_phy_s1_q02_1.png}{170.7}{0}
\qfigdef{p2_phy_s1_q04_1}{Score-II - Test 25 (20 Mar 2024)-img/p2_phy_s1_q04_1.png}{107.2}{0}
\qfigdef{p2_phy_s1_q05_1}{Score-II - Test 25 (20 Mar 2024)-img/p2_phy_s1_q05_1.png}{73.5}{0}
\qfigdef{p2_phy_s2_para7_8_1}{Score-II - Test 25 (20 Mar 2024)-img/p2_phy_s2_para7_8_1.png}{115.8}{0}
\qfigdef{p2_phy_s3_para1_2_1}{Score-II - Test 25 (20 Mar 2024)-img/p2_phy_s3_para1_2_1.png}{143.7}{0}
\qfigdef{p2_phy_s3_para3_4_1}{Score-II - Test 25 (20 Mar 2024)-img/p2_phy_s3_para3_4_1.png}{125.8}{0}
\qfigdef{p2_phy_s4_q07_1}{Score-II - Test 25 (20 Mar 2024)-img/p2_phy_s4_q07_1.png}{138.6}{0}
\qfigdef{p2_phy_s4_q09_1}{Score-II - Test 25 (20 Mar 2024)-img/p2_phy_s4_q09_1.png}{55.6}{0}
\qfigdef{p2_chem_s1_q01_1}{Score-II - Test 25 (20 Mar 2024)-img/p2_chem_s1_q01_1.png}{117.3}{0}
\qfigdef{p2_chem_s1_q01_2}{Score-II - Test 25 (20 Mar 2024)-img/p2_chem_s1_q01_2.png}{111.1}{0}
\qfigdef{p2_chem_s1_q02_3}{Score-II - Test 25 (20 Mar 2024)-img/p2_chem_s1_q02_3.png}{44.8}{0}
\qfigdef{p2_chem_s1_q02_1}{Score-II - Test 25 (20 Mar 2024)-img/p2_chem_s1_q02_1.png}{103.8}{0}
\qfigdef{p2_chem_s1_q02_2}{Score-II - Test 25 (20 Mar 2024)-img/p2_chem_s1_q02_2.png}{100.2}{0}
\qfigdef{p2_chem_s1_q06_1}{Score-II - Test 25 (20 Mar 2024)-img/p2_chem_s1_q06_1.png}{179.0}{0}
\qfigdef{p2_chem_s3_para1_2_1}{Score-II - Test 25 (20 Mar 2024)-img/p2_chem_s3_para1_2_1.png}{310.6}{1}
\qfigdef{p2_chem_s4_q07_1}{Score-II - Test 25 (20 Mar 2024)-img/p2_chem_s4_q07_1.png}{60.1}{0}
% ---- body skeleton ----
\begin{document}
\begin{center}{\Large\bfseries JEE (Main+Advanced) Score-II --- Test 25 (20 Mar 2024)}\end{center}
\begin{multicols}{2}
\raggedcolumns
\papersec{1}
\partsec{1}{PHYSICS}
\sect{Section-I}

\qn{1} The resistance of a wire is $4\Omega$ at 50 $^\circ$C and $5\Omega$ at 100$^\circ$C. The resistance of the wire at 0$^\circ$C will be :
\begin{qoptions}
\opt{A}{$3\Omega$}
\opt{B}{$2\Omega$}
\opt{C}{$1\Omega$}
\opt{D}{$4\Omega$}
\end{qoptions}

\qn{2} Two spherical conductors A and B of radii $R$ and $4R$ respectively, are separated by a large distance. If same charge is given to both the spheres and later they are connected by a conducting wire, then in equilibrium condition, the ratio of the magnitude of the electric fields at the surface of spheres A and B is :-
\begin{qoptions}
\opt{A}{1 : 4}
\opt{B}{4 : 1}
\opt{C}{1 : 2}
\opt{D}{2 : 1}
\end{qoptions}

\qn{3} The pressure at the bottom of a tank of water is $4P$ where $P$ is the atmosphere pressure. If the water is drawn out till level of water is lowered by one-fifth, the pressure at the bottom of the tank will now be
\begin{qoptions}
\opt{A}{$2P$}
\opt{B}{$\dfrac{13}{5}P$}
\opt{C}{$\dfrac{17}{5}P$}
\opt{D}{$\dfrac45P$}
\end{qoptions}

\qn{4} A vernier calliper has 20 divisions on the vernier scale which coincide with 19 divisions on the main scale. The least count of the instrument is 0.1 mm. The length of one main scale division is
\begin{qoptions}
\opt{A}{0.5 mm}
\opt{B}{1 mm}
\opt{C}{2 mm}
\opt{D}{0.25 mm}
\end{qoptions}

\qn{5} A thin prism of angle $4^\circ$ is placed at a distance of 10 cm from the object. The distance of the image from the object is (given $\mu$ of prism $=1.5$)
\begin{qoptions}
\opt{A}{$\dfrac{\pi}{8}$ cm}
\opt{B}{$\dfrac{\pi}{12}$ cm}
\opt{C}{$\dfrac{\pi}{9}$ cm}
\opt{D}{$\dfrac{\pi}{7}$ cm}
\end{qoptions}

\qn{6} When a resistor of $10\Omega$ is connected across a battery, the current is 0.5 A. When a resistor of $5\Omega$ is connected across the same battery, the current is 0.8 A. The emf and internal resistance of the battery are
\begin{qoptions}
\opt{A}{$10\mathrm{V},\ \dfrac{10}{3}\Omega$}
\opt{B}{$20\mathrm{V},\ \dfrac{20}{3}\Omega$}
\opt{C}{$\dfrac{20}{3}\mathrm{V},\ \dfrac{10}{3}\Omega$}
\opt{D}{$\dfrac{40}{3}\mathrm{V},\ \dfrac{20}{3}\Omega$}
\end{qoptions}

\qn{7} In a metallic conductor, under the effect of applied electric field, the free electrons of the conductor :-
\begin{qoptions}
\opt{A}{drift from higher potential to lower potential.}
\opt{B}{move in the curved paths from lower potential to higher potential.}
\opt{C}{move with the uniform velocity throughout from lower potential to higher potential.}
\opt{D}{move in the straight line path in the same direction as applied electric field.}
\end{qoptions}

\qn{8} \textbf{Statement-I :} More the value of magnetic flux linked with a coil, more will be the induced e.m.f. developed in the coil.\\
\textbf{Statement-II :} Lenz's law is direct consequence of law of conservation of energy.
\begin{qoptions}
\opt{A}{Statement-I is false but Statement-II is true.}
\opt{B}{Statement-I is true but Statement-II is false.}
\opt{C}{Both Statement-I and II are false.}
\opt{D}{Both Statement-I and II are true.}
\end{qoptions}

\qn{9} If each diode has a forward bias resistance of $75\Omega$ in the below circuit. Which of the following options is correct?
\qfig{p1_phy_s1_q09_1}
\begin{qoptions}
\opt{A}{$\dfrac{I_2}{I_3}=1$}
\opt{B}{$\dfrac{I_1}{I_2}=2$}
\opt{C}{$\dfrac{I_3}{I_4}=1$}
\opt{D}{$\dfrac{I_1}{I_2}=1$}
\end{qoptions}

\qn{10} A point mass moves along a circle of radius $R$ with a constant angular acceleration $\alpha$. How much time is needed after motion begins for the radial acceleration of the point mass to be equal to its tangential acceleration?
\begin{qoptions}
\opt{A}{$\sqrt{\dfrac1\alpha}$}
\opt{B}{$\sqrt{\dfrac R\alpha}$}
\opt{C}{$\sqrt{\dfrac2\alpha}$}
\opt{D}{$\sqrt{\dfrac1{2\alpha}}$}
\end{qoptions}

\qn{11} A 400 N block is dragged along a horizontal surface by an applied force $\vec F$ as shown. The coefficient of kinetic friction is $\mu_k=0.2$ and the block moves at constant velocity. The magnitude of $\vec F$ is : (close to)
\qfig{p1_phy_s1_q11_5}
\begin{qoptions}
\opt{A}{100 N}
\opt{B}{154 N}
\opt{C}{87 N}
\opt{D}{290 N}
\end{qoptions}

\qn{12} Figure shows a capillary rise $H$. If the air is blown through the horizontal tube in the direction as shown, then rise in capillary tube will become :-
\qfig{p1_phy_s1_q12_1}
\begin{qoptions}
\opt{A}{$=H$}
\opt{B}{$>H$}
\opt{C}{$<H$}
\opt{D}{zero}
\end{qoptions}

\qn{13} A mixture of light, consisting of wavelength 590 nm and an unknown wavelength, illuminates Young's double-slit setup and gives rise to two overlapping interference patterns on the screen. The central maximum of both patterns coincide. Further, it is observed that the third bright fringe of known light coincides with the 4$^{\mathrm{th}}$ bright fringe of the unknown light. From this data, the wavelength of the unknown light is :-
\begin{qoptions}
\opt{A}{393.4 nm}
\opt{B}{885.0 nm}
\opt{C}{442.5 nm}
\opt{D}{776.5 nm}
\end{qoptions}

\qn{14} Consider a uniform wire of mass $M$ and length $L$. It is bent into a semicircle. Its moment of inertia about a line perpendicular to the plane of the wire passing through the centre of semicircle is :-
\begin{qoptions}
\opt{A}{$\dfrac{ML^2}{\pi^2}$}
\opt{B}{$\dfrac14\dfrac{ML^2}{\pi^2}$}
\opt{C}{$\dfrac12\dfrac{ML^2}{\pi^2}$}
\opt{D}{$\dfrac25\dfrac{ML^2}{\pi^2}$}
\end{qoptions}

\qn{15} Position-time graph of particle of mass 2 kg is shown in figure. Total work done on the particle from $t=0$ to $t=4$s is
\qfig{p1_phy_s1_q15_1}
\begin{qoptions}
\opt{A}{8 J}
\opt{B}{4 J}
\opt{C}{0 J}
\opt{D}{can't be determined}
\end{qoptions}

\qn{16} In photo-electric effect, when graph is plotted for stopping potential versus frequency of incident light, for different metals, then :-
\begin{qoptions}
\opt{A}{slope is different for different metals.}
\opt{B}{threshold frequency does not very with emitter.}
\opt{C}{the slope is universal constant equal to $(h/e)$.}
\opt{D}{stopping potential decreases with increasing frequency of incident light.}
\end{qoptions}

\qn{17} A vessel contains an ideal monoatomic gas which expands at constant pressure, when heat 400 J is given to it. Then the work done in expansion is :-
\begin{qoptions}
\opt{A}{400 J}
\opt{B}{240 J}
\opt{C}{160 J}
\opt{D}{80 J}
\end{qoptions}

\qn{18} At a certain height $h$ above the surface of the earth the change in the value of acceleration due to gravity ($g$) is the same as that at a depth $x$ below the surface. Assuming $h$ and $x$ to be enough small compared to the radius of the earth, $x:h$ is
\begin{qoptions}
\opt{A}{1 : 1}
\opt{B}{2 : 1}
\opt{C}{1 : 2}
\opt{D}{1 : 4}
\end{qoptions}

\qn{19} A person is swimming with a speed of 10 m/s at an angle of $120^\circ$ with the flow and reaches to a point directly opposite on the other side of the river. The speed of the flow is `$x$' m/s. The value of `$x$' is.
\begin{qoptions}
\opt{A}{1}
\opt{B}{4}
\opt{C}{8}
\opt{D}{5}
\end{qoptions}

\qn{20} A travelling wave is described by the equation\\
$y(x,t)=0.05\sin(8x-4t)$ m\\
The velocity of the wave is :- [all the quantities are in SI unit]
\begin{qoptions}
\opt{A}{4 ms$^{-1}$}
\opt{B}{2 ms$^{-1}$}
\opt{C}{0.5 ms$^{-1}$}
\opt{D}{8 ms$^{-1}$}
\end{qoptions}

\sect{Section-II}

\qn{1} Two blocks of masses $m_1=5$kg and $m_2=2$ kg are connected by threads which pass over the pulleys as shown in the figure. The threads are massless and the pulleys are massless and smooth. The blocks can move only along the vertical direction. If the acceleration of $m_1$, if expressed in simplest form, is equal to $\left(\dfrac{2n}{n-1}\right)$ m/sec$^2$, calculate `$n$'. (take $g=10$ m/s$^2$)
\qfig{p1_phy_s2_q01_1}

\qn{2} A man in a balloon throws a stone downwards with a speed of 5 ms$^{-1}$ with respect to the balloon which is moving upwards with a constant acceleration of 5 ms$^{-2}$. Find the speed (in ms$^{-1}$) of the stone relative to the man after 2s. (Take $g=10$ ms$^{-2}$)

\qn{3} Solid sphere A is rotating about an axis PQ. If the radius of the sphere is 5 cm then its radius of gyration about PQ will be $\sqrt{x}$ cm. The value of $x$ is \underline{\hspace{3em}}.
\qfig{p1_phy_s2_q03_1}

\qn{4} A body of mass 3 kg collides elastically with another body at rest and then continues to move in the original direction with one-half of its original speed. What is the mass of the target body (in kg)?

\qn{5} A ball of mass 0.45 kg which is initially at rest is hit by a bat. The bat remains in contact with the ball for $3\times10^{-3}$ s. During this time period, the force on the ball by the bat varies with time $t$(in s) as $F(t)=[(\alpha\times10^6)t-(\beta\times10^9)t^2]$ N where $\alpha$ and $\beta$ are constants. The ball's speed, immediately as it loses contact with the bat is 20 m/s. If the relation between $\alpha$ and $\beta$ is $\alpha-2\beta=n$, find $n$.

\qn{6} A hole is drilled in a metal sheet. At 27$^\circ$C, the diameter of hole is 5 cm. When the sheet is heated to 177$^\circ$C, the change in the diameter of hole is $d\times10^{-3}$ cm. The value of $d$ will be \underline{\hspace{3em}} if coefficient of linear expansion of the metal is $1.6\times10^{-5}$/$^\circ$C.

\qn{7} A parallel plate capacitor has plate area 100m$^2$ and plate separation of 10m. The space between the plates is filled upto a thickness 5 m with a material of dielectric constant of 10. The resultant capacitance of the system is `$x$' pF The value of $\varepsilon_0=8.85\times10^{-12}$ F.m$^{-1}$. The value of `$x$' to the nearest integer is \underline{\hspace{2.5em}}

\qn{8} The time period of oscillation of a simple pendulum is given by $T=2\pi\sqrt{\ell/g}$. The length of the pendulum is measured as $\ell=10\pm0.1$cm and the time period as $T=0.5\pm0.02$ s. Determine percentage error in the value of $g$.

\qn{9} In the circuit shown in the figure, the ratio of the quality factor and the band width (in sec) is :-
\qfig{p1_phy_s2_q09_1}

\qn{10} As shown in the figure, a combination of a thin plano concave lens and a thin plano convex lens is used to image an object placed at infinity. The radius of curvature of both the lenses is 30 cm and refraction index of the material for both the lenses is 1.75. Both the lenses are placed at distance of 40 cm from each other. Due to the combination, the image of the object is formed at distance $x=$ \underline{\hspace{2.5em}} cm, from concave lens.
\qfig{p1_phy_s2_q10_1}

\partsec{2}{CHEMISTRY}
\sect{Section-I}

\qn{1} CuSO$_4$ is heated with borax in flame as per given diagram (a) and (b).
\qfig{p1_chem_s1_q01_1}
What observation in bead (after cooling) is obtained in (a) and (b) process respectively.
\begin{qoptions}
\opt{A}{Colorless and Blue-Green}
\opt{B}{Blue-Green and Blue}
\opt{C}{Red opaque and Colorless}
\opt{D}{Blue-Green and Red opaque}
\end{qoptions}

\qn{2} Major product (C) is :
\qfig{p1_chem_s1_q02_1}
\begin{qoptions}
\opt{A}{\qfigi{p1_chem_s1_q02_2}}
\opt{B}{\qfigi{p1_chem_s1_q02_3}}
\opt{C}{\qfigi{p1_chem_s1_q02_4}}
\opt{D}{\qfigi{p1_chem_s1_q02_5}}
\end{qoptions}

\qn{3}\par\nopagebreak
\qfig{p1_chem_s1_q03_1}
Correct statement about C.
\begin{qoptions}
\opt{A}{It gives turbidity in 5 min with lucas reagent}
\opt{B}{It is Ph$-$CHO}
\opt{C}{It is Ph$-$CH$_2-$OH}
\opt{D}{It gives racimisation with PCl$_3$}
\end{qoptions}

\qn{4} 10 gm of \qfigi{p1_chem_s1_q04_1} is dissolved in water to make 100 ml solution. It is kept in 30 cm polarimeter tube. It produced an optical rotation of $-3.6^\circ$. Calculate $[\alpha]$ specific angle of rotation for \qfigi{p1_chem_s1_q04_2} under identical circumstances of temperature.
\begin{qoptions}
\opt{A}{$-12^\circ$}
\opt{B}{$+12^\circ$}
\opt{C}{0}
\opt{D}{$-36^\circ$}
\end{qoptions}

\qn{5} Reaction correctly matched with their major product is/are -
\begin{qoptions}
\opt{A}{$\mathrm{CH_3{-}CH_2{-}COOK}\ \xrightarrow{\mathrm{Electrolysis}}\ \mathrm{CH_3{-}CH_3}$}
\opt{B}{\qfigi{p1_chem_s1_q05_1}}
\opt{C}{\qfigi{p1_chem_s1_q05_3} \qfigi{p1_chem_s1_q05_2}}
\opt{D}{$\mathrm{Me_3C{-}Cl+CH_3O^-\rightarrow Me_3C{-}OCH_3}$}
\end{qoptions}

\qn{6}\par\nopagebreak
\qfig{p1_chem_s1_q06_1}
Choose the incorrect option ?
\begin{qoptions}
\opt{A}{(S) is 4-bromoaniline}
\opt{B}{(Q) is Acetanilide}
\opt{C}{Basic strength order is : (Q) $>$ (P) $>$ (R)}
\opt{D}{Both (B) and (C)}
\end{qoptions}

\qn{7} Which of the reaction will not produce aldehyde as major product :
\begin{qoptions}
\opt{A}{$\mathrm{CH_3{-}C{\equiv}N}\ \xrightarrow[\mathrm{(ii)\ H_2O}]{\mathrm{(i)\ SnCl_2/HCl}}$}
\opt{B}{$\mathrm{CH_3{-}CH_2{-}C{\equiv}N}\ \xrightarrow[\mathrm{(ii)\ H_2O}]{\mathrm{(i)\ DIBAL{-}H}}$}
\opt{C}{\qfigi{p1_chem_s1_q07_1}\ $\xrightarrow{\mathrm{H_2/Pd}}$}
\opt{D}{\qfigi{p1_chem_s1_q07_2}\ $\xrightarrow[\mathrm{Anhyd.\ AlCl_3}]{\mathrm{CO,\ HCl}}$}
\end{qoptions}

\qn{8} In which of the following reaction, one of the product obtained will give isocynide test ?
\begin{qoptions}
\opt{A}{\qfigi{p1_chem_s1_q08_1}\ $\xrightarrow[\mathrm{KOH}]{\mathrm{Br_2}}$}
\opt{B}{\qfigi{p1_chem_s1_q08_2}\ $\xrightarrow{\mathrm{LAH}}$}
\opt{C}{\qfigi{p1_chem_s1_q08_3}\ $\xrightarrow[\mathrm{(ii)\ H_3O^+}]{\mathrm{(i)\ CH_3{-}Br}}$}
\opt{D}{All of these}
\end{qoptions}

\qn{9} In which triad, first one has the highest boiling point.
\begin{qoptions}
\opt{A}{PH$_3$, AsH$_3$, SbH$_3$}
\opt{B}{CH$_3$O$-$O$-$CH$_3$, CH$_3-$S$-$CH$_3$, CH$_3-$Se$-$CH$_3$}
\opt{C}{HBr, HCl, HF}
\opt{D}{AlF$_3$, SiF$_4$, PF$_6$}
\end{qoptions}

\qn{10} Bromine can be liberated from potassium bromide solution by like action of :
\begin{qoptions}
\opt{A}{Iodine solutions}
\opt{B}{Chlorine water}
\opt{C}{Sodium chloride}
\opt{D}{Potassium iodide}
\end{qoptions}

\qn{11} ICl$_2^+$, ICl$_2^-$, ICl$_4^-$. sum of the bond pairs and lone pairs on each iodine atom in the given ionic species :
\begin{qoptions}
\opt{A}{2, 2, 4}
\opt{B}{2, 3, 2}
\opt{C}{4, 5, 6}
\opt{D}{4, 5, 4}
\end{qoptions}

\qn{12} \textbf{Assertion (A) :} Fullerene are the only pure form of carbon.\\
\textbf{Reason (R) :} They have smooth structure without having `dangling' bonds.
\begin{qoptions}
\opt{A}{A is false but R is true.}
\opt{B}{A is true but R is false.}
\opt{C}{Both A and R are true and R is the correct explanation of A.}
\opt{D}{Both A and R are true and R is NOT the correct explanation of A.}
\end{qoptions}

\qn{13}\par\nopagebreak
\newcolumntype{L}[1]{>{\raggedright\arraybackslash}p{#1}}%
\begin{tabular}[t]{|L{0.36\linewidth}|L{0.4\linewidth}|}
\hline
\multicolumn{1}{|c|}{List-I} & \multicolumn{1}{c|}{List-II}\\ \hline
(P) $[\mathrm{CrCl_2(ox)_2}]^{3-}$ & (1) Geometrical isomerism\\ \hline
(Q) $[\mathrm{Co(NH_3)_5(NO_2)}]$ $\mathrm{(NO_3)_2}$ & (2) Ionization isomerism\\ \hline
(R) $[\mathrm{Co(NH_3)_6}]$ $[\mathrm{Co(CN)_6}]$ & (3) Optical isomerism\\ \hline
(S) $[\mathrm{Co(en)_3}]\mathrm{Cl_3}$ & (4) Linkage isomerism\\ \hline
\end{tabular}\\[2pt]
The correct option is :-
\begin{qoptions}
\opt{A}{P $\rightarrow$ 1,3 ; Q $\rightarrow$ 2,4,3 ; R $\rightarrow$ 3,4 ; S $\rightarrow$ 2,3}
\opt{B}{P $\rightarrow$ 1,4 ; Q $\rightarrow$ 2,4 ; R $\rightarrow$ 4 ; S $\rightarrow$ 3}
\opt{C}{P $\rightarrow$ 1,3 ; Q $\rightarrow$ 2,4 ; R $\rightarrow$ 4 ; S $\rightarrow$ 3}
\opt{D}{P $\rightarrow$ 1,3 ; Q $\rightarrow$ 2,4,3 ; R $\rightarrow$ 1,4 ; S $\rightarrow$ 3,2}
\end{qoptions}

\qn{14} An ideal gas undergoes isothermal compression from 15 m$^3$ to 10 m$^3$ against a constant external pressure of 6 Nm$^2$. Heat released in this process is used to increase the temperature of 3 moles of a metal (M), with heat capacity 18 J-mole$^{-1}$ K$^{-1}$, then the temperature of metal increases by :
\begin{qoptions}
\opt{A}{$\dfrac23$ kelvin}
\opt{B}{$\dfrac32$ kevlin}
\opt{C}{$\dfrac45$ kevlin}
\opt{D}{$\dfrac59$ kevlin}
\end{qoptions}

\qn{15} In Balmer series of hydrogen line spectrum, the 4th line from the red end corresponds to which one of the following inter-orbit jump of the electron for Bohr orbits in an atom of hydrogen ?
\begin{qoptions}
\opt{A}{$5\rightarrow2$}
\opt{B}{$3\rightarrow2$}
\opt{C}{$6\rightarrow2$}
\opt{D}{$4\rightarrow2$}
\end{qoptions}

\qn{16} For a reaction $2P+Q\rightarrow$ product, rate law is :\\
$-\dfrac{d[P]}{dt}=k[P]$\\
At a time when $t=\dfrac2k$, concentration of the reactant is `P' : ($C_0$ = initial concentration of reactant `P')
\begin{qoptions}
\opt{A}{$\dfrac{C_0}{e}$}
\opt{B}{$eC_0$}
\opt{C}{$\dfrac{C_0}{e^2}$}
\opt{D}{$\dfrac1{C_0}$}
\end{qoptions}

\qn{17} 0.1 M NaI (aq.) (100\% dissociation) and 0.1 M glucose aqueous solution are mixed in 3 : 2 volume ratio, the depression in freezing point of the resulting solution will be :\\
(Given : $K_f$(H$_2$O) = 1.86 K-kg/mole) (Assume molarity = molality)
\begin{qoptions}
\opt{A}{3.72 K}
\opt{B}{1.86 K}
\opt{C}{0.297 K}
\opt{D}{0.279 K}
\end{qoptions}

\qn{18} For the cell reaction :\\
$2\mathrm{Fe(s)}+\mathrm{O_2(g)}+4\mathrm{H^+(aq)}\longrightarrow2\mathrm{Fe^{2+}(aq)}+2\mathrm{H_2O}(\ell)$\\
at $[\mathrm{Fe^{+2}}]=10^{-4}$ M, $p(\mathrm{O_2})=0.01$ atm and pH = 4, the cell potential at 25$^\circ$C will be :\\
$\left\{\dfrac{2.303RT}{F}=0.06\right\}$ (Given : $E^\circ_{\mathrm{cell\ reaction}}=1.68$ V)
\begin{qoptions}
\opt{A}{1.47 V}
\opt{B}{1.53 V}
\opt{C}{1.77 V}
\opt{D}{1.83 V}
\end{qoptions}

\qn{19} Find solubility (in M) of AgCN in buffer solution of pH = 4. If $K_{sp}$ of AgCN in water $=1.6\times10^{-16}$, $K_a$ of HCN $=6.4\times10^{-10}$
\begin{qoptions}
\opt{A}{$1.6\times10^{-5}$ M}
\opt{B}{$1.9\times10^{-5}$ M}
\opt{C}{$5\times10^{-6}$ M}
\opt{D}{$2.7\times10^{-6}$ M}
\end{qoptions}

\qn{20} \textbf{Statement-I :} There is a change in enthalpy of an ideal gas during compression at constant temperature.\\
\textbf{Statement-II :} Enthalpy of fixed amount of an ideal gas is a function of temperature only.
\begin{qoptions}
\opt{A}{Statement-I is false but Statement-II is true.}
\opt{B}{Statement-I is true but Statement-II is false.}
\opt{C}{Both Statement-I and II are false.}
\opt{D}{Both Statement-I and II are true.}
\end{qoptions}

\sect{Section-II}

\qn{1} Calculate double bond equivalent of compound (M)
\qfig{p1_chem_s2_q01_1}

\qn{2}\par\nopagebreak
\qfig{p1_chem_s2_q02_1}
What is the value of (X + Y)?

\qn{3} Find number of compound which are aromatic in nature :
\qfig{p1_chem_s2_q03_1}

\qn{4} Among $[\mathrm{NiCl_4}]^{-2}$, $\mathrm{Na_3[CoF_6]}$, $\mathrm{CsO_2}$, $\mathrm{Na_2O_2}$, $[\mathrm{Co(NH_3)_4Cl_2}]^+$, $\mathrm{O_2}$, $\mathrm{Ti^{3+}}$\\
Total number of diamagnetic compounds are :

\qn{5} Total number of correct statements are :\\
(i) NO is heavier than O$_2$\\
(ii) The formula of heavy water is D$_2$O\\
(iii) N$_2$ diffuses faster than oxygen through orifice\\
(iv) NH$_3$ can be used as a refrigerant

\qn{6} Sum of bond order of NO , CO , O$_2^-$ , C$_2^{2-}$ is :-

\qn{7} Consider all the elements from Lanthanum ($Z=57$) to Lutetium ($Z=71$) and find number of element(s) which is/are smaller in size as compared to Europium ($Z=63$).

\qn{8} When 2 mole of a non volatile, solid electrolyte is mixed in 20 moles of water, its vapour pressure decreases by 20\%. When 10 gm of same electrolyte is dissolved in 500 gm of water, the solution freezes at $-0.93$ $^\circ$C then find molecular weight of electrolyte (in g/mol).\\
($K_f$(H$_2$O) = 1.86 K-kg/mol and assuming degree of dissociation of electrolyte remains same in both the solutions)

\qn{9} 1 mole of an ideal gas at 400 kelvin in thermal contact with surrounding expands isothermally from 2L to 3L against constant pressure 4 bar, then in this process, calculate magnitude of the change in entropy of surrounding in J/K (1 bar-litre = 100 J)

\qn{10} Calculate pH of 0.01 M CH$_3$NH$_3^+$Cl$^-$ (aq.) solution. [$K_b$(CH$_3$NH$_2$) $=10^{-6}$]

\partsec{3}{MATHEMATICS}
\sect{Section-I}

\qn{1} Exhaustive value of $x$ such that\\
$\cos^{-1}\left(\dfrac{8x}{1+16x^2}\right)=-\dfrac\pi2+2\tan^{-1}(4x)$
\begin{qoptions}
\opt{A}{$\left(\dfrac18,\infty\right)$}
\opt{B}{$\left(\dfrac1{16},\dfrac18\right)$}
\opt{C}{$\left[\dfrac14,\infty\right)$}
\opt{D}{$\left(\dfrac18,\dfrac14\right)$}
\end{qoptions}

\qn{2} If $z$ is a complex number such that $|z|\geq3$ then minimum value of $|z+1|$
\begin{qoptions}
\opt{A}{is equal to 2}
\opt{B}{lies in $(0,2)$}
\opt{C}{is strictly greater than 2}
\opt{D}{lies in $\left(\dfrac52,3\right)$}
\end{qoptions}

\qn{3} \textbf{Assertion (A)} : Number of terms in expansion of $(a+b+c)^{20}$ is 231.\\
\textbf{and}\\
\textbf{Reason (R)} : Total number of non-negative integral solutions of equation $x+y+z=n$ is ${}^{n+2}C_3$.
\begin{qoptions}
\opt{A}{(A) is True, (R) is True ; (R) is a correct explanation for (A).}
\opt{B}{(A) is True, (R) is True ; (R) is not a correct explanation for (A).}
\opt{C}{(A) is True, (R) is False.}
\opt{D}{(A) is False, (R) is True.}
\end{qoptions}

\qn{4} The sum of real roots of the equation $(2x-1)^2-3|2x-1|+2=0$ is -
\begin{qoptions}
\opt{A}{3}
\opt{B}{2}
\opt{C}{1}
\opt{D}{0}
\end{qoptions}

\qn{5} Largest number of the following numbers which divides $\left(2025^{2024}-2021^{2024}\right)$, is :
\begin{qoptions}
\opt{A}{4}
\opt{B}{8}
\opt{C}{16}
\opt{D}{32}
\end{qoptions}

\qn{6} Parabola $y^2=4a(x-c_1)$ \& $x^2=4a(y-c_2)$ touch each other where $c_1$ \& $c_2$ are variables, than locus of their point of contact is :
\begin{qoptions}
\opt{A}{$xy=a^2$}
\opt{B}{$xy=2a^2$}
\opt{C}{$xy=4a^2$}
\opt{D}{None of these}
\end{qoptions}

\qn{7} If $x$ and $y$ are the solutions of the equation $12\sin x+5\cos x=2y^2-8y+21$, then the least positive value of $\dfrac xy$ is -
\begin{qoptions}
\opt{A}{$\cot^{-1}\left(\dfrac5{12}\right)$}
\opt{B}{$\dfrac12\cot^{-1}\left(\dfrac5{12}\right)$}
\opt{C}{$\dfrac14\tan^{-1}\left(\dfrac{12}5\right)$}
\opt{D}{$\dfrac12\cot^{-1}\left(\dfrac{12}5\right)$}
\end{qoptions}

\qn{8} If $A$ is a square matrix of order $n$ then,\\
$\underbrace{(\mathrm{adj\ adj\ adj\ adj}\ldots\mathrm{adj}A)}_{(n-1)\text{times}}\ .\ \underbrace{(\mathrm{adj\ adj\ adj}\ldots\mathrm{adj}A)}_{n\ \text{times}}$\\
is equal to :
\begin{qoptions}
\opt{A}{$\left|A^{n-1}\right|^{(n-1)^2}$}
\opt{B}{$\left|A^{-1}\right|^{(n-1)^{(n-1)}}$}
\opt{C}{$\dfrac1n|A|^{(n-1)^{(n-1)}}\cdot I_{n\times n}$}
\opt{D}{$|A|^{(n-1)^{(n-1)}}\cdot I_{n\times n}$}
\end{qoptions}

\qn{9}\par\nopagebreak
\newcolumntype{L}[1]{>{\raggedright\arraybackslash}p{#1}}%
\begin{tabular}[t]{|L{0.56\linewidth}|c|c|}
\hline
\multicolumn{1}{|c|}{List-I} & & List-II\\ \hline
(A) The vector $\left(\vec a+3\vec b\right)$ is perpendicular to $\left(7\vec a-5\vec b\right)$ and vector $\left(\vec a-4\vec b\right)$ is perpendicular to $\left(7\vec a-2\vec b\right)$. Then angle between $\vec a$ \& $\vec b$ is & (I) & $30^\circ$\\ \hline
(B) For non-zero two vectors $\vec a$ and $\vec b$. If $\left|\vec a\cdot\vec b\right|=\left|\vec a\times\vec b\right|$, then the angle between $\vec a$ \& $\vec b$ is & (II) & $45^\circ$\\ \hline
(C) Angle between the line $\dfrac{x+3}{4}=\dfrac{y-4}{-2}=\dfrac{z+5}{-1}$ and plane $4x-2y-z=1$ is & (III) & $60^\circ$\\ \hline
(D) If $a,b,c$ are different real numbers and $a\hat i+b\hat j+c\hat k$; $b\hat i+c\hat j+a\hat k$ \& $c\hat i+a\hat j+b\hat k$ are position vectors of three non-collinear points A, B \& C then angle between $\overrightarrow{AB}$ and $\overrightarrow{AC}$ is & (IV) & $90^\circ$\\ \hline
\end{tabular}
\begin{qoptions}
\opt{A}{(A)$\rightarrow$(III), (B)$\rightarrow$(II), (C)$\rightarrow$(IV), (D)$\rightarrow$(III)}
\opt{B}{(A)$\rightarrow$(III), (B)$\rightarrow$(IV), (C)$\rightarrow$(II), (D)$\rightarrow$(I)}
\opt{C}{(A)$\rightarrow$(II), (B)$\rightarrow$(III), (C)$\rightarrow$(IV), (D)$\rightarrow$(I)}
\opt{D}{(A)$\rightarrow$(I), (B)$\rightarrow$(III), (C)$\rightarrow$(IV), (D)$\rightarrow$(III)}
\end{qoptions}

\qn{10} A data consists of $n$ observations : $x_1\ x_2\ x_3\ \ldots.x_n$. If $\displaystyle\sum_{i=1}^{n}(x_i+1)^2=9n$ and $\displaystyle\sum_{i=1}^{n}(x_i-1)^2=5n$ the the standard deviation of this data is :
\begin{qoptions}
\opt{A}{7}
\opt{B}{2}
\opt{C}{$\sqrt5$}
\opt{D}{$\sqrt7$}
\end{qoptions}

\qn{11} The number of real values of `$x$' for which\\[4pt]
$\begin{vmatrix}x^3-x-1&2x^3&2x^3\\2x&x-1-x^3&2x\\2&2&1-x-x^3\end{vmatrix}=0$ is :
\begin{qoptions}
\opt{A}{1}
\opt{B}{2}
\opt{C}{18}
\opt{D}{14}
\end{qoptions}

\qn{12} The coordinates of the orthocentre of the triangle formed by the lines $2x^2-2y^2+3xy+3x+y+1=0$ and $3x+2y+1=0$ are -
\begin{qoptions}
\opt{A}{$\left(\dfrac45,\dfrac35\right)$}
\opt{B}{$\left(\dfrac15,-\dfrac45\right)$}
\opt{C}{$\left(\dfrac25,\dfrac15\right)$}
\opt{D}{$\left(-\dfrac35,-\dfrac15\right)$}
\end{qoptions}

\qn{13} The set of all possible values of $p$ for which the point $([p+1],2)$ lying inside the circle $x^2+y^2-x-y-8=0$ is $[a,b)$, where $[.]$ denotes greatest integer function, then $\displaystyle\int_{\frac{a\pi}{8}}^{\frac{b\pi}{8}}\dfrac{x^7-2x^5+x^3-x+1}{\cos^2x}dx$ is equal to -
\begin{qoptions}
\opt{A}{2}
\opt{B}{0}
\opt{C}{4}
\opt{D}{6}
\end{qoptions}

\qn{14} The area bounded by the line $x=1$ and the curve $\sqrt{\dfrac yx}+\sqrt{\dfrac xy}=4$ is :
\begin{qoptions}
\opt{A}{$4\sqrt3$}
\opt{B}{$\sqrt3$}
\opt{C}{$2\sqrt3$}
\opt{D}{$3\sqrt3$}
\end{qoptions}

\qn{15} \textbf{Statement-1 :} If algebraic sum of length of perpendiculars drawn from points A(4, 2), B(2, 4) and C(0, 0) to variable line is zero, then this line passes through (2, 2).\\
\textbf{and}\\
\textbf{Statement-2 :} If algebraic sum of length of perpendiculars drawn from three non-collinear points A($x_1$, $y_1$), B($x_2$, $y_2$) and C($x_3$, $y_3$) to a variable line is zero, then this line passes through centroid of $\Delta$ABC.
\begin{qoptions}
\opt{A}{Both Statement-I and Statement-II is true}
\opt{B}{Both Statement-I and Statement-II is false}
\opt{C}{Statement-I is true but Statement-II is false}
\opt{D}{Statement-I is false but Statement-II is true}
\end{qoptions}

\qn{16} If coefficient of $x^2y^3z$ in $(x+y+z)^n$ is $A$ (where $A\ne0$) then coefficient of $xy^4z$ is :
\begin{qoptions}
\opt{A}{$2A$}
\opt{B}{$\dfrac{nA}{2}$}
\opt{C}{$\dfrac A2$}
\opt{D}{$nA$}
\end{qoptions}

\qn{17} Two dice are thrown simultaneously to get the co-ordinates on $x$-$y$ plane. Then the probability that this point lies inside or on the region bounded by $|x|+|y|=3$, is :
\begin{qoptions}
\opt{A}{$\dfrac3{14}$}
\opt{B}{$\dfrac23$}
\opt{C}{$\dfrac1{12}$}
\opt{D}{$\dfrac4{14}$}
\end{qoptions}

\qn{18} If $S$ represent the sum of the solutions of the equation $\sin\theta=\sin\left(\tan^{-1}\cot\left(\dfrac{75\pi}{8}\right)\right)$ in $[0,2\pi]$, then $\cos S$ is -
\begin{qoptions}
\opt{A}{$\dfrac12\sqrt{2+\sqrt2}$}
\opt{B}{$\dfrac12\sqrt{2-\sqrt2}$}
\opt{C}{1}
\opt{D}{$-1$}
\end{qoptions}

\qn{19} Consider the lines $L_1$ \& $L_2$ given by\\
$L_1:\dfrac{x-1}2=\dfrac{y-3}1=\dfrac{z-2}2$ \&\\
$L_2:\dfrac{x-2}1=\dfrac{y-2}2=\dfrac{z-3}3$\\
A line having direction ratio $1,-1,-2$ intersects $L_1$ \& $L_2$ at the points P \& Q respectively. Find the length of PQ.
\begin{qoptions}
\opt{A}{$2\sqrt6$}
\opt{B}{$3\sqrt2$}
\opt{C}{$4\sqrt3$}
\opt{D}{4}
\end{qoptions}

\qn{20} If the ellipse $\dfrac{x^2}{16}+\dfrac{y^2}{b^2}=1$ is inscribed in a square of side length $5\sqrt2$, then area of ellipse is
\begin{qoptions}
\opt{A}{$4\pi$}
\opt{B}{$8\pi$}
\opt{C}{$10\pi$}
\opt{D}{$12\pi$}
\end{qoptions}

\sect{Section-II}

\qn{1} In a sequence of positive integers, each term after the second is the product of the previous two terms. If sixth term is 4000, then the sum of first two terms

\qn{2} If $f:R\rightarrow R$ is a function such that $f(x)=x^3+x^2f'(1)+f''(2)x+f'''(4)\ \forall\ x\in R$, then find $f(1)$

\qn{3} The number of points of discontinuity of $f(x)$ in $[0,2]$ where\\
$f(x)=\begin{cases}[\cos\pi x]&x\leqslant1\\|2x-3|[x-2]&x>1\end{cases}$\\
( [.] denotes the greatest integer function)

\qn{4} The curve $x+y-\ln(x+y)=2x+5$ has a vertical tangent at the point $(\alpha,\beta)$. Then $\alpha+\beta$ is equal to :

\qn{5} If $I_n=\displaystyle\int_0^\pi e^x(\sin x)^ndx$ then $10\,\dfrac{I_3}{I_1}$ is equal to :

\qn{6} Solution of differential equation\\
$x^3\dfrac{dy}{dx}=y^3+y^2\sqrt{y^2-x^2}$\\
is $cxy=y+k\sqrt{y^2-x^2}$ where $c$ is arbitrary constant if $k$ is equal to :

\qn{7} If the circle $(x-a)^2+y^2=25$ intersects the circle $x^2+(y-b)^2=16$ in such a way that common chord is of maximum length, then the value of $a^2+b^2$ is equal to \underline{\hspace{3em}}

\qn{8} Let $3f(x)+2f\left(\dfrac{x+59}{x-1}\right)=10x+30$; $x\ne1$, then the value of $f(7)$ is

\qn{9} Entries of a $2\times2$ determinant are chosen from the set $\{-1,1\}$. The probability that determinant has zero value is equal to $P$ then $4P$ is \underline{\hspace{2em}}.

\qn{10} If the absolute maximum value of $f(x)=\dfrac1{|x-4|+1}+\dfrac1{|x+8|+1}$ is $M$, then $\dfrac{13M}7$ is equal to \underline{\hspace{2em}}.

\papersec{2}
\partsec{1}{PHYSICS}
\sect{Section-I(i)}

\qn{1} If $(N+2)$ equal mass hang from a system of pulley (all pulleys are ideal and upward direction is taken to be positive), then :-
\qfig{p2_phy_s1_q01_1}
\begin{qoptions}
\opt{A}{for $N=5$ acceleration of end mass is $\dfrac{-5g}{11}$ m/s$^2$}
\opt{B}{for $N=3$ acceleration of end mass is $\dfrac{-3g}{7}$ m/s$^2$}
\opt{C}{for $N=5$ acceleration of all the mass except end mass is $\dfrac{g}{11}$ m/s$^2$}
\opt{D}{for $N=3$ acceleration of all the mass except end mass is $\dfrac{g}{7}$ m/s$^2$}
\end{qoptions}

\qn{2} A free particle of mass $m$ is confined to a region of length $L$. The de-Broglie wave associated with the particle is sinusoidal in nature as given in the figure then choose the correct options.
\qfig{p2_phy_s1_q02_1}
\begin{qoptions}
\opt{A}{Momentum of particle is $p=\dfrac hL$.}
\opt{B}{Momentum of particle is $p=\dfrac{2h}L$.}
\opt{C}{Energy of particle is $E=\dfrac{2h^2}{mL^2}$}
\opt{D}{Energy of particle is $E=\dfrac{h^2}{2mL^2}$}
\end{qoptions}

\qn{3} A gas mixture contains $n_1$ mole of gas 1 (molar mass $M_1$) and $n_2$ mole of gas 2 (molar mass $M_2$) with adiabatic ratio $\gamma_1$ and $\gamma_2$ respectively with molar specific heats $C_{p_1}$, $C_{v_1}$ and $C_{p_2}$, $C_{v_2}$. Mixture is kept at temperature $T$ (in K) then choose the correct option.
\begin{qoptions}
\opt{A}{Adiabatic ratio of mixture is $\gamma_{mix}=\dfrac{n_1\gamma_1+n_2\gamma_2}{n_1+n_2}$.}
\opt{B}{Adiabatic ratio of mixture is $\gamma_{mix}=\dfrac{n_1C_{p_1}+n_2C_{p_2}}{n_1C_{v_1}+n_2C_{v_2}}$.}
\opt{C}{rms speed of molecule in mixture $v_{rms}=\sqrt{\dfrac{3RT}{\left(\frac{n_1M_1+n_2M_2}{n_1+n_2}\right)}}$.}
\opt{D}{rms speed of molecule in mixture $v_{rms}=\sqrt{\dfrac{3RT}{\left(\frac{n_1+n_2}{\frac{n_1}{M_1}+\frac{n_2}{M_2}}\right)}}$.}
\end{qoptions}

\qn{4} A liquid of density `$\rho$' is filled in a cuboidal container of dimension ($H\times H\times2H$) as shown in the figure. A semicircular gate is made at the bottom of container of radius $R=\dfrac H2$. Let $F$ is the force acting on the gate due to the liquid and $\tau$ is the torque acting on the gate due to the liquid about a line PQ. Then choose the correct option(s).
\qfig{p2_phy_s1_q04_1}
\begin{qoptions}
\opt{A}{$F=\dfrac{H^3\rho g}{4}\left(\dfrac\pi4-\dfrac13\right)$}
\opt{B}{$F=\dfrac{H^3\rho g}{2}\left(\dfrac\pi4-\dfrac16\right)$}
\opt{C}{$\tau=\dfrac{H^4\rho g}{4}\left(\dfrac13-\dfrac\pi{32}\right)$}
\opt{D}{$\tau=\dfrac{H^3\rho g}{3}\left(\dfrac13-\dfrac\pi{16}\right)$}
\end{qoptions}

\qn{5} Capacitor of capacitance $C=10\mu$F was connected with a source of emf $\varepsilon=10$ V for long time. At time $t=0$ switch S is shifted from point A to B, which connects a solenoid of self inductance $L=10$H, length $\ell=\dfrac12$m, and number of turns 10000. A charge particle of charge $q=10\mu$C, mass $1\times10^{-6}$ kg is released from rest on the axis of solenoid at time $t=\dfrac{10\pi}{3}$ms. At this instant :-
\qfig{p2_phy_s1_q05_1}
\begin{qoptions}
\opt{A}{acceleration of the charged particle is zero.}
\opt{B}{acceleration of the charged particle is 100 m/s$^2$.}
\opt{C}{magnetic field inside the solenoid is nearly 700 mT.}
\opt{D}{magnetic field inside the solenoid is nearly 218 $\mu$T.}
\end{qoptions}

\qn{6} Which of the following statement(s) is/are CORRECT ?
\begin{qoptions}
\opt{A}{If the energy of bombarding electron in the Coolidge tube is increased then the cutoff wavelength decreases while the characteristic wavelength remains unaltered.}
\opt{B}{If the energy of bombarding electron in the Coolidge tube is increased then the cutoff wavelength increases while the characteristic wavelength remains unaltered.}
\opt{C}{The characteristic wavelength depends on the target metal while the cutoff wavelength is a function of the energy of bombarding electrons.}
\opt{D}{The characteristic wavelength is function of the energy of the bombarding electrons while the cutoff wavelength depends on the characteristic of target metal.}
\end{qoptions}

\sect{Section-I(ii)}

\paraq{7}{8}
If $n$ balls, $B_1,B_2,B_3,\ldots..B_n$ having mass $m_1,m_2,m_3,\ldots,m_n$ with ($m_1\gg m_2\gg m_3\gg\ldots>m_n$) standing in a vertical stake. The bottom of $B_1$ is at height `$h$' above the ground and the bottom of $B_n$ is a height ($h+\ell$) above the ground. (All collisions are elastic)
\qfig{p2_phy_s2_para7_8_1}

\qn{7} If balls are dropped then to what height does the top ball bounce? If number of balls is 4 and $\ell$ is 10 m and height is 2m.
\begin{qoptions}
\opt{A}{460}
\opt{B}{230}
\opt{C}{115}
\opt{D}{320}
\end{qoptions}

\qn{8} If $h=1.8$ m then find minimum number of ball so that the top ball escape from earth :-\\
(Assume $v_{escape}=11400$ m/s)
\begin{qoptions}
\opt{A}{9}
\opt{B}{10}
\opt{C}{11}
\opt{D}{12}
\end{qoptions}

\paraq{9}{10}
Acoustic impedance of a medium is defined as the force per unit layer velocity that a given region applies to an adjacent region, as a wave propagates. Remember that impedance is a property of the medium and not the wave, even though it is generally easiest to calculate it by considering the properties of a travelling wave.

Let us consider a longitudinal progressive wave $s=s_0\sin(\omega t-kx)$.\\
The velocity of a "layer" of molecules whose equilibrium position is $x$ is given by
\[V_P=\frac{\partial s}{\partial t}\]
To find the force, consider a cross-sectional area $A$. We can write the (excess) force that the sheet exerts on the region to its right as
\[F=A\Delta P=-A\left(-B\frac{ds}{dx}\right)\]
And since we are working with a travelling wave we can write.
\[F=\mp A\frac Bv\frac{ds}{dt}\]
where $v$ is velocity of wave and $B$ is bulk modulus of medium.

The force that the sheet feels from the region on its right is the negative of this, but the sign isn't important when calculating the impedance $Z$, because $Z$ is defined to be the magnitude of $\dfrac F{V_P}$, so impedance
\[Z=\frac F{V_P}=\frac{(A)(B)}v=\frac{A\rho v^2}v=A\rho v\]
Where $\rho$ is density of medium.

The impedance per unit area is the more natural thing to talk about, because $Z/A$ is independent of the specific cross section chosen. The force $F$ can be made arbitrarily large by making the area $A$ arbitrarily large, so $Z=F/v$ isn't too meaningful. When people talk about the impedance of air, they usually mean "impedance per area", that is, force per velocity per area. When a sound wave gets reflected from a boundary while going from medium 1 to 2, the amplitudes of reflected and transmitted wave are given by
\[\Delta P_r=\frac{Z_2-Z_1}{Z_2+Z_1}\Delta P_i\]
\[\Delta P_t=\frac{2Z_2}{Z_1+Z_2}\Delta P_i\]
where $\Delta P_i$ is amplitude of incident wave.
A longitudinal wave goes from medium 1 to medium 2 where $\rho_1v_1=5\times10^6$ kg/m$^2$-s and $\rho_2v_2=45\times10^6$ kg/m$^2$-s, where $\rho$ and $v$ refers to density of medium and wave velocity in the medium.

\qn{9} Find ratio of pressure amplitude of transmitted wave to pressure amplitude of incident wave.
\begin{qoptions}
\opt{A}{1.25}
\opt{B}{1.80}
\opt{C}{0.80}
\opt{D}{1.50}
\end{qoptions}

\qn{10} Find ratio of displacement amplitude of transmitted to displacement amplitude of incident wave.
\begin{qoptions}
\opt{A}{5.00}
\opt{B}{0.50}
\opt{C}{0.20}
\opt{D}{2.00}
\end{qoptions}

\sect{Section-II(i)}

\paraq{1}{2}
A cylindrical vessel filled with water upto height $H=1$ m is kept fixed on an inclined plane having inclination angle $\theta$ as shown in figure. A small hole is made in the vessel at a height `$h$' from base of vessel so that water stream cover a range $R$ on the inclined plane.
\qfig{p2_phy_s3_para1_2_1}

\qn{1} If $\theta=15^\circ$, find $h$ (in m) for which range on incline plane $R$ will be maximum.

\qn{2} If $\theta=30^\circ$, find the maximum range (in m) on incline plane covered by water stream.

\paraq{3}{4}
A charged particle of mass $m$ and charge $q$ is connected to a inextensible string tied to the floor and projected on a rough horizontal $xy$-plane surface with $z$-axis in the vertically upward direction. Both electric and magnetic fields are acting in the region and given by $\vec E=-E_0\hat k$ and $\vec B=-B_0\hat k$ respectively. The particle enters into the field at $(a,0,0)$ with velocity $\vec v=v_0\hat j$. The particle starts moving into a circular path on the plane. If the coefficient of friction between the particle and the plane is $\mu$. Then calculate the (Consider $m=1$ kg, $q=2$C, $v_0=10$ m/s, $E_0=5$ N/C, $\mu=0.4$, $g=10$ m/s$^2$) :
\qfig{p2_phy_s3_para3_4_1}

\qn{3} Time when the particle will come to rest (in s) is :-

\qn{4} Distance travelled by the particle when it comes to rest (in m) is :-

\paraq{5}{6}
A wire was laid along the equator on the supports. The wire is made of aluminum with a cross section area of 1 cm$^2$. The magnetic field near the equator is homogeneous, horizontal, directed to the north and equal in magnitude to $3\times10^{-5}$ T. Initially temperature is room temperature $=20^\circ$ C. Density of aluminum: $D=2.7$ g cm$^3$, melting point: $t=660^\circ$C, specific heat of fusion: $L=390$ kJ/kg. Resistivity of aluminum is $2.7\times10^{-8}\ \Omega$m at the room temperature and remains constant. Specific heat capacity of aluminum is 0.90 J/g $^\circ$C. The dimension of the wire also remain same.

\qn{5} What current (in $10^4$ A) do I need to pass through the wire so that it ceases to put pressure on the supports (begins to levitate)?

\qn{6} Estimate the minimum acceptable duration of this experiment (in sec), if the wire does not have time to lose heat into the surrounding space :-

\sect{Section-II(ii)}

\qn{7} A thin convex lens having focal length 20 cm is cut into two parts (A and B) 10 mm above the principle axis. The lower part `A is placed with its optical centre at origin and upper part `B' at (90 cm, 0) as shown in the figure. A point object `P' is placed at ($-30$ cm, 0). If the distance of final image from $x$-axis is $10\times N$ mm, then find the value of $N$. (Assuming paraxial ray approximation to remain valid)
\qfig{p2_phy_s4_q07_1}

\qn{8} On the axis of a long tube with a mirror inner surface are an isotropic point source of light and a completely absorbing incident light ball with a radius of 12 cm. The center of the ball is separated by distance 15 cm from the source. What should be maximum internal diameter (in cm) of the pipe so that the ball absorbs exactly half of the energy emitted by the source?

\qn{9} A vertical infinite line charge has linear charge density $\lambda=4\pi\epsilon_0$ SI units. A charge of $q=40\times10^{-6}$ C and mass $3\times10^{-6}$ kg is suspended from the line as shown. Find the horizontal distance of charge to the line (in cm). The length of thread is 1m. ($g=10$ m/s$^2$) (Round off to nearest integer).
\qfig{p2_phy_s4_q09_1}

\partsec{2}{CHEMISTRY}
\sect{Section-I(i)}

\qn{1} Choose the correct statement
\begin{qoptions}
\opt{A}{Dipole moment of \qfigi{p2_chem_s1_q01_1}}
\opt{B}{Resonance energy of \qfigi{p2_chem_s1_q01_2}}
\opt{C}{$\mathrm{pK_{a_2}}$ of Maleic acid $>$ $\mathrm{pK_{a_2}}$ of Fumaric acid}
\opt{D}{Basic strength of Et$_2$NH $>$ nBu $-$ NH$_2$}
\end{qoptions}

\qn{2} Reactions which can give benzene as a product :
\begin{qoptions}
\opt{A}{\qfigi{p2_chem_s1_q02_3}\ $\xrightarrow[\Delta]{\mathrm{Al_2O_3\,/\,Cr_2O_3}}$}
\opt{B}{$\mathrm{CH{\equiv}CH}\ \xrightarrow{\mathrm{Red\ hot\ Fe-Cu\ tube}}$}
\opt{C}{\qfigi{p2_chem_s1_q02_1}}
\opt{D}{\qfigi{p2_chem_s1_q02_2}}
\end{qoptions}

\qn{3} Select correct statement about thiazyl dichloride ion $\left(\mathrm{NSCl_2^-}\right)$ and thionyl chloride $(\mathrm{OSCl_2})$
\begin{qoptions}
\opt{A}{NSCl$_2^-$ is isoelectronic with OSCl$_2$}
\opt{B}{OSCl$_2$ has smaller Cl--S--Cl angle than NSCl$_2^-$}
\opt{C}{NSCl$_2^-$ has longer S$-$Cl bond than OSCl$_2$}
\opt{D}{Both contains $\mathrm{p_\pi-d_\pi}$ bond}
\end{qoptions}

\qn{4} $\mathrm{HNO_3+P_2O_5\rightarrow HPO_3+X}$\\
$\mathrm{X+I_2\rightarrow Y+Z}$\\
$\mathrm{Y+CO\rightarrow I_2+CO_2}$\\
$\mathrm{X+NaCl\rightarrow P+Q}$\\
(P shows brown ring test)\\
Which of the following option(s) is/are correct
\begin{qoptions}
\opt{A}{X is acidic oxide}
\opt{B}{Z is colourless gas}
\opt{C}{Y is used for estimation of CO}
\opt{D}{Hybridisation of central atom in Q is sp$^2$}
\end{qoptions}

\qn{5} SO$_3$ gas is entering the environment in Kota at a constant rate of $6.93\times10^{-6}$ gm/L/day due to the emission of polluting gases from thermal power plant, but at the same time it is decomposing \& following first order kinetics with half life of 100 days.\\
(Assume SO$_3$ present in Kota with constant concentration means rate of formation and dissociation of SO$_3$ are equal)\\
Based on above information select the true statement(s).
\begin{qoptions}
\opt{A}{Concentration of SO$_3$ in Kota is $1.25\times10^{-5}$ M}
\opt{B}{If SO$_3$ emission is stopped then after 1000 days its concentrations will reduce to $1.2\times10^{-8}$ M.}
\opt{C}{If $10^3$ L of air is passed through 1 L pure water (assuming all SO$_3$ to be dissolved in it) \& resulting solution is titrated against 1 N, NaOH solution of which 15 ml is consumed.}
\opt{D}{An industry is manufacturing H$_2$SO$_4$ at the rate of 980 kg per day with the use of SO$_3$ in air it should use $8\times10^5$ Liter air/day.}
\end{qoptions}

\qn{6}\par\nopagebreak
\qfig{p2_chem_s1_q06_1}
\noindent\textit{Each beaker:} 1 litre solution containing C$_2$H$_5$OH in 1 M Na$_2$S, 1 M NaHS\\
0.5 F of charge is used for the electrolysis of above solution. Ethyl alcohol is oxidised to enthanal at anode and water is reduced at cathode.\\
Information $\Rightarrow$ $\left(\mathrm{P^{Ka_1}}\right)$, $\left(\mathrm{P^{Ka_2}}\right)$ values of H$_2$S are $x$, $y$ respectively.\\
(I) Volume of solution is constant during electrolysis.\\
(II) Na$_2$S, NaHS do not participate in electrolytic reactions\\
(III) Gram molar volume at STP is 22.7 litres.\\
Which of the following is/are correct?
\begin{qoptions}
\opt{A}{pH of anode after electrolysis is $y-\log3$}
\opt{B}{pH of cathode after electrolysis is $y+\log3$}
\opt{C}{0.25 mole ethanal is produced at anode due to electrolysis}
\opt{D}{Volume of H$_2$ produced at cathode at STP is 5.675 litres.}
\end{qoptions}

\sect{Section-I(ii)}

\paraq{7}{8}
Pseudo halides like CN$^-$, N$_3^-$, SCN$^-$ can exhibit properties similar to halide ion such as Cl$^-$, Br$^-$, I$^-$, in solubility \& complex formation.

\qn{7} The compound which is soluble in water at room temperature is :
\begin{qoptions}
\opt{A}{AgSCN}
\opt{B}{NH$_4$SCN}
\opt{C}{AgCN}
\opt{D}{Pb(CN)$_2$}
\end{qoptions}

\qn{8} $\mathrm{CN^-}\ \xrightarrow[\mathrm{ion(X)}]{\mathrm{Metal}}\ \mathrm{C_2N_2}$\\
$\mathrm{SCN^-}\ \xrightarrow[\mathrm{ion(Y)}]{\mathrm{Metal}}\ \mathrm{S_2C_2N_2}$\\
X and Y are respectively :-
\begin{qoptions}
\opt{A}{Hg$^{2+}$, Fe$^{+3}$}
\opt{B}{Cu$^{+2}$, Fe$^{+2}$}
\opt{C}{Fe$^{+3}$, Cu$^{+2}$}
\opt{D}{Cu$^{+2}$, Cu$^{+2}$}
\end{qoptions}

\paraq{9}{10}
Adsorption is the accumulation of molecular species at the surface. During an experiment CH$_3$COOH molecules were adsorbed on the charcoal having surface area $\sqrt3\times10^3$m$^2$. Assume CH$_3$COOH molecule is spherical with radius 1\AA\ \& whole available charcoal surface is occupied by CH$_3$COOH molecules forming monolayer in two dimensional close packing with maximum efficiency ($N_A=6\times10^{23}$)

\qn{9} What will be number of moles of CH$_3$COOH adsorbed on charcoal -
\begin{qoptions}
\opt{A}{1/6}
\opt{B}{1/12}
\opt{C}{2/3}
\opt{D}{1/3}
\end{qoptions}

\qn{10} If after adsorption molarity reduced to 1/3M for CH$_3$COOH solution. Find \%( w/v) for 500ml solution used initially -
\begin{qoptions}
\opt{A}{0.5}
\opt{B}{3}
\opt{C}{4}
\opt{D}{6}
\end{qoptions}

\qfullbegin
\sect{Section-II(i)}

\paraq{1}{2}
One of the best methods for the regioselective alkylation of ketones consists in an activation of the required $\alpha$-carbon atom. For example :
\qfig{p2_chem_s3_para1_2_1}
\qfullend

\qn{1} Degree of unsaturation of `D$_2$' is-

\qn{2} Number of carbon in `D$_1$'

\paraq{3}{4}
An anion `A' do not gives volatile product with dilute H$_2$SO$_4$ but gives brown ring test is passed through Zn + NaOH evolving gas `X'.

\qn{3} When gas `X' is passed through a chemical substance Y under basic medium yields brown precipitate of $\left[\mathrm{Hg_2(NH_2)(O)}\right]^+\mathrm{I^-}$. The EAN of the metal in `Y' is.

\qn{4} When gas `X' is passed through green solution of NiCl$_2$ a deep blue colored solution is observed due to the formation of complex `Z'.\\
What is the sum of primary valency \& secondary valency of `Z'.

\paraq{5}{6}
Decomposition of phosgene gas, $\mathrm{COCl_2}\ \overset{k_f}{\underset{k_b}{\rightleftharpoons}}\ \mathrm{CO+Cl_2}$ follows first order kinetics with :\\
$\log k_f=13.48-\dfrac{52400}{4.5}\left(\dfrac1T\right)$\\[0pt]
[$K_f$, $K_b$ are rate constants of forward and backward reaction respectively at $T$K]\\
$\log k_b=8.023-\dfrac{46600}{4.5}\left(\dfrac1T\right)$

\qn{5} Determine $\Delta_rH$ (in kJ) for the decomposition of one mole of phosgene gas.

\qn{6} Determine $\Delta_rS$ (in J/K) for decomposition reaction.

\sect{Section-II(ii)}

\qn{7} Identify total number of stereoisomers of following compound :
\qfig{p2_chem_s4_q07_1}

\qn{8} Identify total number of Essential amino acids\\
(a) Glycine \quad (b) Alanine\\
(c) Phenyl alanine \quad (d) Arginine\\
(e) Cysteine \quad (f) Histidine\\
(g) Proline \quad (h) Tryptophan

\qn{9} Total number of metal cations given below which can be detected by H$_2$S gas in acidic medium?\\
Cd$^{2+}$, Pb$^{2+}$, Cu$^{2+}$, Hg$^{2+}$, Co$^{2+}$, Zn$^{2+}$, Ca$^{2+}$, Ba$^{2+}$

\partsec{3}{MATHEMATICS}
\sect{Section-I(i)}

\qn{1} If $f$ is a differentiable function such that $f(x+y)=f(x)f(y)\ \forall\ x,y\in R$, $f(0)\ne0$ and $f'(0)=1$. Then
\begin{qoptions}
\opt{A}{$\displaystyle\int_{-1}^{1}\dfrac{(\ell n(f(x)))^2}{1+(f(x))^2}dx=\dfrac13$}
\opt{B}{$\displaystyle\int_{-\pi/2}^{\pi/2}\dfrac{\cos(\ell n(f(x)))}{1+f(x)}dx=1$}
\opt{C}{$\displaystyle\int_{-2024}^{2024}\ell n\left(\ell n(f(x))+\sqrt{1+(\ell n(f(x)))^2}\right)dx=0$}
\opt{D}{$\displaystyle\int_0^2(f(x)-\ell n(f(x)))dx=e^2-2$}
\end{qoptions}

\qn{2} $A=\begin{bmatrix}0&1&2\\0&-3&0\\1&1&1\end{bmatrix}$ and $6A^{-1}=aA^2+bA+cI$ where $a,b,c\in$ integer and $I$ is identity matrix then
\begin{qoptions}
\opt{A}{$a+b+c=-2$}
\opt{B}{$A^{-1}=\dfrac16\begin{bmatrix}-3&1&6\\0&-2&0\\3&1&0\end{bmatrix}$}
\opt{C}{$a+b+c=18$}
\opt{D}{$A^{-1}=\dfrac1{12}\begin{bmatrix}3&3&2\\2&4&-1\\5&1&-4\end{bmatrix}$}
\end{qoptions}

\qn{3} For any two events $A$ and $B$ of sample space $S$, if $P(A)=\dfrac35$ and $P(B)=\dfrac23$ then which of the following options are correct ?
\begin{qoptions}
\opt{A}{$P\left(A\cap\overline B\right)\leqslant\dfrac13$}
\opt{B}{$P(A\cup B)\geqslant\dfrac23$}
\opt{C}{$\dfrac4{15}\leqslant P(A\cap B)\leqslant\dfrac35$}
\opt{D}{$\dfrac25\leqslant P(A\cap B)\leqslant\dfrac9{10}$}
\end{qoptions}

\qn{4} For a $3\times3$ invertible matrix $A$ satisfying the characteristic equation $A^3+pA^2+qA-rI=0$, which of the following is/are true [Here $\mathrm{tr}(A)$ = trace of matrix $A$, $\det(A)$ = determinant value of matrix $A$]
\begin{qoptions}
\opt{A}{$p=-\mathrm{tr}(A)$}
\opt{B}{$q=\dfrac{(\mathrm{tr}(A))^2-\mathrm{tr}\left(A^2\right)}{2}$}
\opt{C}{$r=\det(A)$}
\opt{D}{$q=\dfrac{\mathrm{tr}\left(A^2\right)-(\mathrm{tr}(A))^2}{2}$}
\end{qoptions}

\qn{5} If $A_i$ represents the area bounded by $|x-a_i|+|y|=b_i$ ; $i\in$ natural number and $a_{i+1}=a_i+\dfrac32b_i$ and $b_{i+1}=\dfrac{b_i}2$ ; $a_1=0$ ; $b_1=32$ then
\begin{qoptions}
\opt{A}{$A_3=128$}
\opt{B}{$\displaystyle\lim_{n\to\infty}\sum_{i=1}^{n}A_i=\dfrac43(16)^2$}
\opt{C}{$\displaystyle\lim_{n\to\infty}\sum_{i=1}^{n}A_i=\dfrac83(32)^2$}
\opt{D}{$A_3=256$}
\end{qoptions}

\qn{6} If $xf^2(x)-(x^2-x+2)f(x)=x(x-2)$; $\forall\ x\in R$ and $f(2)\ne0$ then $\displaystyle\int_0^2\left(\dfrac{f(x)+f(-x)+f''(x)}{2-f'(x)}\right)f(x)dx$ is equal to
\begin{qoptions}
\opt{A}{$\displaystyle\int_{-1}^{1}f(x)dx$}
\opt{B}{$\displaystyle\int_0^2f(x)dx$}
\opt{C}{$\displaystyle\int_0^1(f(x)+f(-x))\,dx$}
\opt{D}{$2\displaystyle\int_0^1f(x)dx$}
\end{qoptions}

\sect{Section-I(ii)}

\paraq{7}{8}
Let $\alpha$, $\beta$ and $\gamma$ are roots of the equation $x^3+6x+3=0$ and\\
$A=\cos^{-1}\left(\sin\left((\alpha+\beta)^{-1}+(\beta+\gamma)^{-1}+(\gamma+\alpha)^{-1}\right)\right)$\\
$B=\cos\left(\tan^{-1}\left(\sin\left(\dfrac{\alpha+\beta+\gamma}2\right)\right)\right)$\\
$C=\sec^{-1}\left(\mathrm{cosec}(1-\alpha)(1-\beta)(1-\gamma)\right)$

\qn{7} The value of $5A+B-C$
\begin{qoptions}
\opt{A}{21}
\opt{B}{1}
\opt{C}{0}
\opt{D}{None of these}
\end{qoptions}

\qn{8} Range of the function $g(x)=\dfrac{(5A-C)x^5+6Bx^2}{x^4+(B-1)x^3+1}$
\begin{qoptions}
\opt{A}{$[0,3]$}
\opt{B}{$(0,3)$}
\opt{C}{$[0,\infty)$}
\opt{D}{$(-\infty,\infty)$}
\end{qoptions}

\paraq{9}{10}
Given five points on a circle of radius 4 units. The centres of rectangular hyperbolas passing through these points taken four at a time are lying on the fixed curve $S$, then

\qn{9} The curve $S$ is
\begin{qoptions}
\opt{A}{a hyperbola}
\opt{B}{a parabola}
\opt{C}{an ellipse}
\opt{D}{a circle}
\end{qoptions}

\qn{10} Choose the correct option :
\begin{qoptions}
\opt{A}{eccentricity of $S$ is $\sqrt2$}
\opt{B}{any pair of perpendicular tangents of the curve $S$ meet on a line}
\opt{C}{area of $S$ is $4\pi$}
\opt{D}{eccentricity of $S$ is $\dfrac12$}
\end{qoptions}

\sect{Section-II(i)}

\paraq{1}{2}
Let $f(x)$ be a continuous function satisfying $\displaystyle\int_0^1f(x)(4x^2-f(x))dx=\dfrac45$ and $\displaystyle\int\dfrac{16+12f(x)}{72-96\sqrt3x+54(f(x))^2}dx=\dfrac1{6\sqrt6}\tan^{-1}(\alpha x+\beta)-\dfrac1{6\sqrt3g(x)}+C$ where `$C$' is constant of inegration $\left(x\ne\dfrac1{\sqrt3},\ g(0)=-1\right)$ then

\qn{1} The value of $f(3)$ \underline{\hspace{3em}}

\qn{2} The value of $g\left(\sqrt3\right)$ \underline{\hspace{3em}}

\paraq{3}{4}
We are given the curves $y=\displaystyle\int_{-\infty}^{x}f(t)\,dt$ through the point $\left(0,\dfrac12\right)$ and $y=f(x)$ where $f(x)>0$ and $f(x)$ is differentiable $\forall\ x\in R$ through $(0,1)$. If tangents drawn to both the curves at the points with equal abscissae intersect on the same point on the $x$-axis then

\qn{3} Number of solution of $f(x)=2ex$ is \underline{\hspace{3em}}

\qn{4} $\displaystyle\lim_{x\to\infty}(f(x))^{f(-x)}$ \underline{\hspace{3em}}

\paraq{5}{6}
If $P$ is a point on the curve $5x^2+3xy+y^2=2$ and $O$ is the origin, then

\qn{5} Minimum value of $11(OP)^2$ is

\qn{6} Maximum value of $OP^2$ is

\sect{Section-II(ii)}

\qn{7} Three pairs of real numbers $(x_1,y_1)$ $(x_2,y_2)$ and $(x_3,y_3)$ satisfy both equations $x^3-3xy^2=2005$ and $y^3-3x^2y=2004$ find value of $\dfrac{y_1y_2y_3}{2(y_1-x_1)(y_2-x_2)(y_3-x_3)}$

\qn{8} If the total number of non-decreasing functions defined from $f:\{1,2,3,4,5\}\longrightarrow\{1,2,3,4,5,6,7,8,9\}$ is $m$ then $\dfrac m{143}$ is equal to

\qn{9} $x,y$ are independent of each other, also $x,y\in R$ then maximum value of\\
$\left(2\cos^2x+2\sqrt2\sin x\cos x+\sin^2x\right)+$\\
$\dfrac1{2\cos^2y+2\sqrt2\sin y\cdot\cos y+\sin^2y+1}$ is


\end{multicols}
\end{document}
